%% ======================================================================
%% chap8/discussion_main.tex
%%
%% Chapter 8: Discussion
%%
%% Interprets the findings of Chapter 7: answers each of the five
%% research sub-questions, interprets the most significant results
%% (26.2 TPS figure, the audience-size economic crossover at ~3,600
%% followers, the Snowbridge demonstration), revisits the major
%% architectural decisions in light of ecosystem events during the
%% implementation window, situates the work in context, acknowledges
%% five principal limitations, draws implications for four stakeholder
%% groups (designers, creators, regulators, researchers), and offers a
%% critical self-reflection.
%%
%% Sections:
%%   8.1 Answering the Research Questions (5 subsections, one per SQ)
%%   8.2 Interpretation of Key Results (3 subsections)
%%   8.3 Architectural Decisions Revisited
%%   8.4 The Work in Context
%%   8.5 Limitations
%%   8.6 Implications
%%   8.7 Critical Self-Reflection
%%   8.8 Summary
%%
%% External labels referenced:
%%   chap:results (Ch 7), chap:conclusion (Ch 9),
%%   sec:technical-results (7.1), sec:economic-results (7.2),
%%   sec:pallet-impl (5.2),
%%   subsec:comparative-analysis (7.1.7),
%%   subsec:decentralization-hhi (7.3.3),
%%   subsec:security-audit (7.3.1),
%%   sec:rmrk-to-pallet-nfts (Appendix A.1),
%%   sec:ink-discontinuation (Appendix A.9),
%%   subsec:bounded-vec-future (9.3.4 in Table 9.1),
%%   subsec:audience-validation-future (9.3.5 in Table 9.1)
%%
%% Citations: madapatiPradhan2025, dengEtAl2025, liEtAl2025a,
%% maricEtAl2025, raoEtAl2024, cirelloEtAl2023, cherdakovEtAl2023,
%% mendozaTelloEtAl2024, lahiriEtAl2026
%% ======================================================================

\chapter{Discussion}\label{chap:discussion}

This chapter analyzes the findings of Chapter~\ref{chap:results} by identifying the questions addressed, the decisions substantiated, and the limitations that shape the interpretation of the contribution.

\section{Answering the Research Questions}\label{sec:answering-rqs}

\subsection{SQ1: Cross-Chain Rights Operations and Finality}\label{subsec:sq1-xcm-rights}

The key finding is that the payer--beneficiary decoupling pattern, combined with XCM's \texttt{Transact} instruction, is sufficient to delegate rights issuance to the authoritative chain without requiring two-phase commit. The remote parachain pays; CCRMS issues the rights record; no distributed rollback is needed because rights state is always authoritative on one chain.

Affirmative with one caveat. XCM v3 is sufficient for executing cross-chain rights operations: the canonical sequence \texttt{WithdrawAsset}~$\to$ \texttt{BuyExecution}~$\to$ \texttt{Transact} covers all five \texttt{xcm\_*} extrinsics and achieves operations that are atomic within each chain's execution; XCM \texttt{Transact} does not guarantee atomicity across chains (Appendix~\ref{app:security}, Finding~A). The destination-chain inclusion delay was measured 0 to 7 blocks (median 5) (Section~\ref{sec:technical-results}). No custom XCM extensions are required. The only caveat is that cross-chain recurring renewals depend on the readiness of XCM~v5 \texttt{Schedule} production; in the meantime, an on-chain block hook manages local auto-renewal. This demonstration is inter-parachain (local HRMP). It is not an end-to-end Ethereum~$\to$~CCRMS rights path (Section~\ref{subsec:non-claims}).

\subsection[SQ2: Latency versus Decentralization]{SQ2: Latency versus Decentralization Trade-Offs}\label{subsec:sq2-latency-decentralization}

The trade-off is real but highly asymmetric: CCRMS operational latency is approximately $67{,}000\times$ higher than the centralized baseline (Table~\ref{tab:comparative}: $\sim$6,000~ms vs.\ 0.09~ms), while a production parachain would inherit an HHI of $\sim$16.7 compared with 10,000 for a centralized monopoly (Table~\ref{tab:hhi}), a decentralization gain of roughly two orders of magnitude on that index, for safety, not for test-network liveness (Section~\ref{subsec:decentralization-hhi}). For discrete, low-frequency rights operations (the CCRMS use case), the trade-off is advantageous; for high-frequency interactive applications it is not.

\subsection[SQ3: Unified Rights Record]{SQ3: Unified Rights Record Enforcing All Three Monetization Models}\label{subsec:sq3-unified-rights-token}

Yes. The unified primitive works (Section~\ref{sec:pallet-impl}): three independent storage maps with a priority-resolution access check enforce all three access semantics without novel cryptographic mechanisms, closing gap~G7. Story Protocol's PIL is a generic licensing primitive configured per policy; CCRMS is a specialized access-mode primitive (one Content record, three maps, one ordered check, and renewal/expiry as storage transitions). That is the claimed distinction; we do not claim PIL cannot be configured for similar commercial outcomes. Royalty propagation applies uniformly to all primary-market revenue. Cost savings are substantial, but audience-dependent (Section~\ref{sec:economic-results}; Section~\ref{subsec:economic-crossover}).

\subsection{SQ4: Security and Authorization}\label{subsec:sq4-security-authorization}

\texttt{xcm\_transfer\_ownership} accepted a \texttt{from} identity parameter without verifying that the caller held it; a one-line \texttt{ensure!} check (Finding~B, Section~\ref{subsec:security-audit}) remediated the issue. The lesson generalizes: \textbf{cross-chain extrinsics that accept identity parameters require explicit authorization checks}.

\subsection[SQ5: Economic Viability]{SQ5: Economic Viability under Realistic Creator Audience Distributions}\label{subsec:sq5-economic-viability}

\textbf{Yes, conditionally.} The Monte Carlo simulation (Section~\ref{sec:economic-results}) indicates, at the level of the design, that CCRMS yields higher mean creator revenue across all four simulated platform configurations, with an advantage ranging from 8.9\% to 16.9\%. The result is robust to the fee range (advantage persists for any CCRMS fee below 17\%; Appendix~\ref{app:monte-carlo}) but depends on the assumed audience distribution. This advantage depends heavily on the size of the creator's audience. We explore the crossover point and its implications in Section~\ref{subsec:economic-crossover}. Our contribution is to transform this conditionality into a \textbf{quantitative and reproducible} framework rather than merely qualitative observations.

\section{Interpretation of Key Results}\label{sec:interpretation}

\subsection{The 26.2 TPS Figure in Context}\label{subsec:26-tps-context}

The figure of 26.2~TPS reflects a single-collator local Zombienet configuration that does not exercise the throughput features of Polkadot~2.0 (Agile Coretime, Asynchronous Backing, and Elastic Scaling were available in the \texttt{stable2512} SDK but not enabled in the chain specification), and does not represent the architecture's maximum capacity. Polkadot achieves scalability through \textbf{horizontal} scaling: with 100 parachains at 26 TPS each, the total capacity reaches 2,600 TPS \citep{polkadotParachainSlots, wood2016}. The 26.2 TPS figure is consistent with the range reported for public blockchain systems at comparable scale \citep{raoEtAl2024}, and horizontal scaling across multiple chains is the principal path to higher aggregate throughput. More importantly, the $270\times$ ratio relative to the centralized baseline is misleading. CCRMS is designed to process low-frequency events; it records discrete rights transactions with cryptographic finality, automatic royalty propagation, and cross-chain portability, none of which the centralized baseline provides at any throughput level. When assessed against its actual workload, 26.2 TPS comfortably exceeds the demands of every creator-economy platform.

\subsection{The Economic Crossover at $\sim$3,600 Followers}\label{subsec:economic-crossover}

The economic advantage of CCRMS primarily benefits creators with established audiences. For those with fewer than $\sim$3,600 followers (with the transition beginning near 2,000 followers where the win rate reaches 6.6\%), centralized platforms tend to generate greater revenue because algorithmic discovery helps creators reach audiences they might not otherwise reach. While this initially appears to challenge the notion that CCRMS serves as a creator-friendly alternative, the more accurate and credible assertion is conditional: (decentralized = better for creators \textbf{with established audiences}). CCRMS is primarily a tool designed \textbf{for creators who have surpassed the discovery benefits of centralized platforms and seek to maximize revenue retention}.


\section{Architectural Decisions Revisited}\label{sec:architectural-decisions-revisited}

\textbf{The pallet-first decision proved well-founded.} The pivot was driven by three engineering constraints (Appendix~\ref{sec:contract-first-pivot}), and it was taken against the backdrop of the January~2026 announcement that ink! development would be discontinued, which removed any expectation that the contract layer would mature. The episode illustrates that architectural decisions reducing dependence on any single non-essential component become more valuable as the ecosystem evolves \citep{parnas1972, bassEtAl2021}. This principle is also reflected in the choice of \texttt{pallet-nfts} over the frozen RMRK pallets and in HRMP rather than custom inter-parachain protocols. The same pattern holds for the Snowbridge integration: the configuration required for Ethereum bridging was completed at minimal additional cost, though the Ethereum-to-CCRMS path itself remains untested.

The unified rights record and the payer--beneficiary pattern are likewise vindicated: the former exhibits linear storage growth (fifty-holder cap noted in Section~\ref{sec:limitations}), and the latter is sound subject to Finding~B's authorization requirement.

\section{The Work in Context}\label{sec:work-in-context}

The Polkadot ecosystem evolved during implementation. \textbf{Polkadot~2.0} features reached production by late 2025, indicating that our latency metrics are pre-2.0 upper bounds; deployment uses the updated SDK with no pallet-level changes. \textbf{Snowbridge V2} is backward-compatible with the V1 prototype. \textbf{DATA Foundation} (formerly Story Protocol, rebranded June~2026 to focus on AI training-data provenance) strengthens the case that the cross-chain creator-rights niche remains unserved; comparisons with the Programmable IP License remain relevant. \textbf{EU regulation} (AI Act, MiCA) drives demand for cryptographic provenance; MiCA's treatment of subscription tokens remains legally unclear. Cross-chain interoperability surveys \citep{dengEtAl2025, liEtAl2025a, maricEtAl2025} provide context but do not change the conclusions.

\section{Limitations}\label{sec:limitations}

The CCRMS prototype has six principal limitations:

\begin{enumerate}
\item \textbf{No production deployment.} All measurements are from a local Zombienet topology with two relay-chain validators, single-collator parachains, and no real-world latency. Chapter~\ref{chap:results} results are upper bounds in this test environment and do not predict production behavior.
\item \textbf{Single-author validation.} A single author performed the implementation, test design, measurement collection, and interpretation. No independent code review, third-party security audit, or external benchmark validation has been conducted. Production deployment would require all three before any non-trivial value passes through the system.
\item \textbf{Local-only Snowbridge integration.} Geth and Lodestar ran on the same machine as the Polkadot side, avoiding many realities of public Ethereum (variable block times, MEV, network contention, beacon chain reorgs). The integration serves as a proof of architectural feasibility but provides no data on production bridge behavior.
\item \textbf{The fifty-holder bound.} The \texttt{BoundedVec<\_, ConstU32<50>{}>} bound on \texttt{Children} limits each content item to fifty child NFTs. Because each first-time purchase mints one for the buyer and expired subscriptions never release theirs, a content item can serve up to fifty subscribers and owners over its lifetime. This places the audience tiers that benefit most (Section~\ref{sec:economic-results}; Section~\ref{subsec:prototype-bound}) beyond the current instantiation. Finding~I is the low-cost exploit of that bound: \texttt{purchase\_views(num\_views = 0)} mints a child at zero price (Table~\ref{tab:security-findings}). The bound is not intrinsic to the design; Appendix~\ref{sec:e-child-cap} describes the mitigation paths, none of which are yet implemented. The Monte Carlo result therefore characterizes the intended design, not the bounded prototype (Claim~3, Section~\ref{subsec:claims}).
\item \textbf{No empirical proof of the bridged-payment final hop.} The bridged-ETH-pays-for-CCRMS-operation flow is conceptually straightforward but has not been tested end-to-end. We claim forward-compatibility with the flow, which is a weaker claim than an empirical demonstration would provide.
\item \textbf{Simulation input sensitivity.} Simulated outcomes are robust to pricing and fee parameters but depend strongly on the assumed audience distribution and value of algorithmic discovery, neither validated against real data; consumer onboarding friction is unmodeled.
\end{enumerate}

\section{Implications}\label{sec:implications}

The findings extend beyond the immediate prototype in four directions.

\textbf{For platform designers,} the overarching lesson is that \textbf{architectural choices that minimize reliance on any single non-essential component pay dividends during ecosystem-level changes}. Two corollaries: \textbf{cross-chain extrinsics that accept identity parameters require explicit authorization audits} (Finding~B), and \textbf{reimplementing an unmaintained standard is preferable to depending on it} (the RMRK 2.0 pivot).

\textbf{For creators,} the Monte Carlo simulation suggests that the question ``Is a decentralized platform actually better for me?'' has a conditional answer: \textbf{yes, depending on audience size}. Creators with over 10,000 followers substantially benefit from lower fees; those with fewer than $\sim$2,000 are better served by centralized discovery. The choice should reflect audience size.

\textbf{For regulators,} on-chain rights management systems should be treated as \textbf{compliance infrastructure} for emerging copyright and AI training regimes rather than as alternative payment platforms. The open question about MiCA's classification of subscription tokens with fungible characteristics would benefit from regulatory clarity, as the current ambiguity poses risks for would-be EU deployers.

\textbf{For the research community,} the prototype, implementation specifications, and Monte Carlo simulation are intended to support follow-up research. The methodological contribution is a worked DSR example for multi-component blockchain prototyping. The empirical contribution is the audience-dependent finding, which future research could refine using real-world migration data or models that incorporate creator psychology and switching costs.

\section{Critical Self-Reflection}\label{sec:self-reflection}

In retrospect, we would have \textbf{initiated the Snowbridge integration earlier}; the setup complexity (Appendix~\ref{app:challenges}) required several days of unexpected debugging. We would also have \textbf{prioritized FRAME benchmarking earlier} to report fee-adjusted throughput. We would \textbf{not} change the pallet-first decision, the unified rights record model, or the cross-chain operation pattern.

\section{Summary}\label{sec:chapter-8-summary}

The five research sub-questions are answered affirmatively; the most economically significant finding is the audience-dependent CCRMS advantage. The major architectural decisions are vindicated, most emphatically the pallet-first decision. The six principal limitations constrain how the contribution should be read but do not invalidate it. Chapter~\ref{chap:conclusion} presents the conclusion and future work.
